\documentclass[a4paper]{article}

% Basic packages
% Español (México) con XeLaTeX: fontspec para fuentes Unicode, polyglossia
% para la división silábica y los nombres del documento en la variante mexicana.
\usepackage{fontspec}
\usepackage{polyglossia}
\setdefaultlanguage[variant=mexican]{spanish}
% Los nombres chinos van como «pinyin (original)»: xeCJK da a los caracteres
% Han su propia fuente; sin ella XeLaTeX los omite con un aviso.
% FandolSong no cubre algunos caracteres de nombres (U+73FA); WenQuanYi Zen Hei sí.
\usepackage[AutoFallBack=true]{xeCJK}
\setCJKmainfont{FandolSong-Regular.otf}
\setCJKfallbackfamilyfont{\CJKrmdefault}{WenQuanYi Zen Hei}
% polyglossia no traduce los nombres de listings.
\AtBeginDocument{\providecommand{\lstlistingname}{}\renewcommand{\lstlistingname}{Listado}%
  \providecommand{\lstlistlistingname}{}\renewcommand{\lstlistlistingname}{Índice de listados}}
\usepackage{amsmath, amssymb}
\usepackage{graphicx}
\usepackage[margin=2.5cm]{geometry}
\usepackage[most]{tcolorbox}
\usepackage{etoolbox}
\usepackage{listings}
\usepackage{hyperref}
\usepackage{booktabs}
\usepackage{subcaption}
\usepackage{float}
\usepackage{tikz}
\IfFileExists{pgfplots.sty}{
    \usepackage{pgfplots}
    \pgfplotsset{compat=1.18}
}{}

% Highlight boxes
\newtcolorbox{knowledgebox}[1]{
    enhanced, colback=blue!5!white, colframe=blue!75!black, colbacktitle=blue!75!black,
    coltitle=white, fonttitle=\bfseries, title={#1}, attach boxed title to top left={yshift=-2mm, xshift=2mm},
    boxrule=1pt, sharp corners
}

\newtcolorbox{importantbox}[1]{
    enhanced, colback=yellow!10!white, colframe=yellow!80!black, colbacktitle=yellow!80!black,
    coltitle=black, fonttitle=\bfseries, title={#1}, sharp corners
}

\newtcolorbox{warningbox}[1]{
    enhanced, colback=red!5!white, colframe=red!75!black, colbacktitle=red!75!black,
    coltitle=white, fonttitle=\bfseries, title={#1}, sharp corners
}

% Listings
\lstset{
    language=Python,
    basicstyle=\ttfamily\small,
    keywordstyle=\color{blue},
    stringstyle=\color{red!60!black},
    commentstyle=\color{green!60!black},
    breaklines=true,
    frame=single,
    numbers=left,
    numberstyle=\tiny\color{gray},
    captionpos=b,
    extendedchars=true
}

% Front-page metadata
\newcommand{\notetitle}{MIT 6.S191: Introduction to Deep Learning}
\newcommand{\noteauthors}{Elaboradas a partir de la clase de Alexander Amini}
\newcommand{\notedate}{Primavera de 2025}
\newcommand{\videochannel}{Alexander Amini (MIT)}
\newcommand{\videopublishdate}{2025-03-03}
\newcommand{\videoduration}{69 minutos}
\newcommand{\videourl}{https://www.youtube.com/watch?v=alfdI7S6wCY}
\newcommand{\videocoverpath}{cover.jpg}
\newcommand{\repourl}{https://github.com/hqhq1025/ai-course-notes}

% Footer with repo info
\usepackage{fancyhdr}
\pagestyle{fancy}
\fancyhf{}
\fancyfoot[C]{\thepage}
\fancyfoot[R]{\footnotesize\color{gray}\href{\repourl}{AI Course Notes}}
\renewcommand{\headrulewidth}{0pt}
\renewcommand{\footrulewidth}{0pt}

\begin{document}

\begin{titlepage}
\centering
{\Large Notas del curso\par}
\vspace{1.2cm}
{\huge\bfseries \notetitle\par}
\vspace{0.8cm}
{\large \noteauthors\par}
\vspace{0.3cm}
{\large \notedate\par}
\vspace{1.2cm}

\ifdefempty{\videocoverpath}{
{\small Al generar las notas, indique la ruta local de la portada del video.\par}
}{
\includegraphics[width=0.82\textwidth,height=0.45\textheight,keepaspectratio]{\videocoverpath}\par
}

\vfill
\begin{tcolorbox}[width=0.9\textwidth, colback=black!2!white, colframe=black!60, sharp corners]
\textbf{Autor o canal del video}: \videochannel\par
\textbf{Fecha de publicación}: \videopublishdate\par
\textbf{Duración del video}: \videoduration\par
\ifdefempty{\videourl}{}{
\textbf{Enlace del video}: \href{\videourl}{\nolinkurl{\videourl}}\par
}
\textbf{Sitio del curso}: \href{https://introtodeeplearning.com}{\nolinkurl{introtodeeplearning.com}}\par
\textbf{Página del proyecto}: \href{\repourl}{\nolinkurl{\repourl}}\par
\end{tcolorbox}
\end{titlepage}

\tableofcontents
\newpage

%% --- Inicio del contenido --- %%

\section{Panorama del curso y avances del aprendizaje profundo}

MIT 6.S191 es el curso introductorio oficial de MIT sobre aprendizaje profundo, impartido por Alexander Amini y Ava Amini. Se trata de un campamento intensivo (boot camp) de una semana de duración que cubre la teoría y la práctica fundamentales del aprendizaje profundo. 2025 es el octavo año en que se imparte el curso, que ha enseñado ya a más de 13 millones de estudiantes en todo el mundo.

\subsection{Una década de evolución del aprendizaje profundo}

El profesor Amini abrió la clase con una comparación elocuente: hace apenas una década (2015), el sistema de generación de rostros más avanzado en aprendizaje profundo solo producía bloques de píxeles borrosos (Goodfellow et al.); para 2018 (Karras, Laine, Aila), los rostros generados ya tenían un realismo fotográfico; en 2020, el equipo del curso MIT 6.S191 llegó incluso a producir un video deepfake en el que una figura de «Obama» daba la bienvenida a los estudiantes, video que acumuló más de 1.1 millones de reproducciones en YouTube.

\begin{figure}[H]
\centering
\includegraphics[width=0.85\textwidth]{slides-images/slide-002.jpg}
\caption{Una década de avances en la generación de rostros mediante aprendizaje profundo: de los píxeles borrosos de 2015 al video realista de 2020\protect\footnotemark}
\end{figure}
\footnotetext{Diapositiva 2.}

\begin{knowledgebox}{El costo de producción del deepfake de 2020}
Producir aquel video deepfake de 2 minutos en 2020 requirió: 2 horas de grabación de audio profesional, 50 horas de datos de video en alta definición, un guion estático predefinido y más de 15,000 dólares en recursos de cómputo. Para 2025, en cambio, Amini hizo una demostración en vivo de clonación de voz en tiempo real: bastaron unos minutos de audio de su clase grabados para generar de inmediato un clon de su voz, con el que sostuvo una conversación dinámica y sin guion. Esto ilustra con claridad el enorme salto que ha dado la IA generativa en tan solo unos años.
\end{knowledgebox}

\begin{figure}[H]
\centering
\includegraphics[width=0.85\textwidth]{slides-images/slide-006.jpg}
\caption{2020 contra 2025: comparación de las capacidades de la IA generativa\protect\footnotemark}
\end{figure}
\footnotetext{Diapositiva 6.}

\subsection{¿Qué es el aprendizaje profundo?}

Antes de entender el aprendizaje profundo, es necesario aclarar tres niveles de conceptos:

\begin{importantbox}{La jerarquía AI $\supset$ ML $\supset$ DL}
\begin{itemize}
\item \textbf{Artificial Intelligence (AI)}: cualquier técnica que permite a una computadora imitar el comportamiento humano
\item \textbf{Machine Learning (ML)}: un método en el que, en lugar de programar explícitamente, se deja que la máquina aprenda patrones a partir de los datos
\item \textbf{Deep Learning (DL)}: un subconjunto de ML que usa redes neuronales profundas para extraer patrones a partir de datos crudos
\end{itemize}
Idea central: \textbf{enseñar a la computadora a aprender directamente de los datos crudos cómo realizar una tarea} (learn a task directly from raw data).
\end{importantbox}

\begin{figure}[H]
\centering
\includegraphics[width=0.85\textwidth]{slides-images/slide-007.jpg}
\caption{La relación anidada entre AI, Machine Learning y Deep Learning\protect\footnotemark}
\end{figure}
\footnotetext{Diapositiva 7.}

La esencia de \textbf{Intelligence} (inteligencia) es la capacidad de procesar información para guiar decisiones futuras. Artificial Intelligence consiste en construir algoritmos artificiales para llevar a cabo ese mismo proceso: usar datos para impulsar decisiones. Machine Learning no le indica directamente a la computadora cómo procesar los datos, sino que deja que descubra patrones en ellos de forma autónoma. Deep Learning, por su parte, lleva este proceso un paso más allá al usar redes neuronales profundas.

\subsection{Estructura del curso}

MIT 6.S191 dura una semana e incluye los siguientes módulos:

\begin{table}[H]
\centering
\begin{tabular}{lll}
\toprule
\textbf{Fecha} & \textbf{Tema de la clase} & \textbf{Laboratorio} \\
\midrule
Day 1 & Introduction to Deep Learning & Lab 1: Generación de música \\
Day 2 & Deep Sequence Modeling & Lab 2: Visión por computadora \\
Day 3 & Deep Computer Vision & Lab 3: Modelos de lenguaje de gran escala \\
Day 4 & Deep Generative Modeling & -- \\
Day 5 & Deep Reinforcement Learning & Final Project \\
\bottomrule
\end{tabular}
\caption{Panorama del calendario del curso}
\end{table}

Entre los puntos destacados de este año se encuentran: laboratorios compatibles simultáneamente con los marcos de trabajo TensorFlow y PyTorch, así como un laboratorio de LLM completamente nuevo (ajuste fino de un modelo Gemma de 2,000 millones de parámetros y construcción de un evaluador de IA).

\subsection{Resumen de la sección}

El aprendizaje profundo ha logrado avances asombrosos en la última década: de la generación de píxeles borrosos a la clonación de voz en tiempo real y a la conversación dinámica. El curso MIT 6.S191 busca, mediante una semana de entrenamiento intensivo, que los estudiantes dominen las técnicas fundamentales que impulsan estos avances. Todo el curso gira en torno a una idea central: enseñar a la computadora a aprender tareas directamente a partir de datos crudos.
\section{¿Por qué el aprendizaje profundo? ¿Por qué ahora?}

\subsection{Las limitaciones del aprendizaje automático tradicional}

El aprendizaje automático tradicional depende de \textbf{características diseñadas a mano} (hand-engineered features). Tomando como ejemplo la detección de rostros, el ingeniero tiene que definir manualmente cómo extraer de la imagen características como bordes, curvas y rasgos faciales. Este método presenta tres problemas fundamentales:

\begin{enumerate}
\item \textbf{Consume tiempo}: la ingeniería de características requiere una gran cantidad de tiempo y experiencia de especialistas del dominio
\item \textbf{Es frágil}: las características diseñadas a mano no son robustas ante cambios del entorno (iluminación, ángulo, etc.)
\item \textbf{No es escalable}: cada tarea nueva exige diseñar de nuevo el conjunto de características
\end{enumerate}

\begin{importantbox}{La idea central del aprendizaje profundo}
El avance clave del aprendizaje profundo consiste en: \textbf{dejar que el modelo aprenda automáticamente, a partir de los datos, representaciones jerárquicas de las características} (hierarchical feature representations). Las características de bajo nivel (como bordes y líneas) se combinan automáticamente en características de nivel medio (como ojos, nariz), y estas a su vez se combinan en características de alto nivel (como la estructura facial). Sin intervención humana.
\end{importantbox}

\begin{figure}[H]
\centering
\includegraphics[width=0.85\textwidth]{slides-images/slide-018.jpg}
\caption{Aprendizaje jerárquico de características: de características de bajo nivel a características de alto nivel\protect\footnotemark}
\end{figure}
\footnotetext{Diapositivas, página 18.}

\subsection{Las tres fuerzas impulsoras}

Aunque la teoría de las redes neuronales se remonta a hace varias décadas (el descenso de gradiente estocástico de 1952, el perceptrón de 1958, la retropropagación de 1986, las redes convolucionales profundas de 1995), su verdadero auge se debe a la maduración simultánea de tres fuerzas impulsoras:

\begin{figure}[H]
\centering
\includegraphics[width=0.85\textwidth]{slides-images/slide-019.jpg}
\caption{La historia de las redes neuronales y las tres fuerzas impulsoras del auge del aprendizaje profundo\protect\footnotemark}
\end{figure}
\footnotetext{Diapositivas, página 19.}

\begin{table}[H]
\centering
\begin{tabular}{lp{10cm}}
\toprule
\textbf{Fuerza impulsora} & \textbf{Manifestación concreta} \\
\midrule
Big Data & La aparición de conjuntos de datos a gran escala (como ImageNet, Wikipedia) y la fuerte caída del costo de recolección y almacenamiento de datos \\
Hardware & La capacidad de cómputo masivamente paralela de las GPU hizo posible entrenar redes profundas \\
Software & Los marcos de código abierto (TensorFlow, PyTorch, Keras) simplificaron la construcción y el entrenamiento de modelos de aprendizaje profundo \\
\bottomrule
\end{tabular}
\caption{Las tres fuerzas impulsoras del auge del aprendizaje profundo}
\end{table}

\begin{knowledgebox}{Perspectiva histórica: estas técnicas no son nuevas}
Amini enfatiza en particular que casi todas las técnicas presentadas en esta clase (el perceptrón, la retropropagación, el descenso de gradiente) se inventaron hace \textbf{varias décadas}. Lo verdaderamente nuevo es la capacidad asombrosa que muestran cuando confluyen big data, un poder de cómputo considerable y buenas herramientas. Comprender estos fundamentos es esencial para seguir haciendo avanzar el campo.
\end{knowledgebox}

\subsection{Resumen de la sección}

La ventaja central del aprendizaje profundo frente al aprendizaje automático tradicional consiste en aprender automáticamente representaciones de características, sin necesidad de ingeniería manual de características. Aunque la teoría de base tiene ya varias décadas de historia, la triple fuerza impulsora de big data, hardware de GPU y software de código abierto ha hecho que el aprendizaje profundo viva un desarrollo explosivo en los últimos años.

\section{El perceptrón: el bloque de construcción básico del aprendizaje profundo}

\subsection{Propagación hacia adelante del perceptrón}

El \textbf{perceptrón} (Perceptron) es la unidad estructural básica de todas las redes neuronales, también llamada «neurona» (neuron). Su proceso de propagación hacia adelante consta de tres pasos:

\begin{enumerate}
\item \textbf{Suma ponderada}: multiplicar cada entrada $x_i$ por su peso correspondiente $w_i$ y luego sumar
\item \textbf{Suma del sesgo}: sumar el término de sesgo $w_0$ (bias)
\item \textbf{Activación no lineal}: obtener la salida mediante una función de activación no lineal $g(\cdot)$
\end{enumerate}

\begin{figure}[H]
\centering
\includegraphics[width=0.85\textwidth]{slides-images/slide-022.jpg}
\caption{Propagación hacia adelante del perceptrón (con término de sesgo)\protect\footnotemark}
\end{figure}
\footnotetext{Diapositivas, página 22.}

Su expresión matemática es:

$$\hat{y} = g\left(w_0 + \sum_{i=1}^{m} x_i w_i\right) = g(w_0 + \mathbf{X}^T \mathbf{W})$$

donde:
\begin{itemize}
\item $\mathbf{X} = [x_1, x_2, \ldots, x_m]^T$: vector de entrada
\item $\mathbf{W} = [w_1, w_2, \ldots, w_m]^T$: vector de pesos
\item $w_0$: término de sesgo (bias), que permite desplazar la función de activación a lo largo del eje $z$
\item $g(\cdot)$: función de activación no lineal
\item $\hat{y}$: salida predicha
\end{itemize}

\begin{importantbox}{La función del término de sesgo}
El término de sesgo $w_0$ es un parámetro indispensable en el perceptrón. Permite desplazar la función de activación con independencia de los valores de entrada, de modo que la neurona pueda producir una salida distinta de cero incluso cuando todas las entradas son cero. Un perceptrón sin sesgo solo puede aprender fronteras de decisión que pasan por el origen.
\end{importantbox}

\subsection{Funciones de activación}

La función central de la función de activación $g(\cdot)$ es \textbf{introducir no linealidad en la red}. Entre las funciones de activación más comunes están:

\begin{figure}[H]
\centering
\includegraphics[width=0.95\textwidth]{slides-images/slide-025.jpg}
\caption{Tres funciones de activación comunes: Sigmoid, Tanh y ReLU\protect\footnotemark}
\end{figure}
\footnotetext{Diapositivas, página 25.}

\begin{table}[H]
\centering
\begin{tabular}{llll}
\toprule
\textbf{Función de activación} & \textbf{Fórmula} & \textbf{Rango de salida} & \textbf{Características} \\
\midrule
Sigmoid & $\sigma(z) = \frac{1}{1+e^{-z}}$ & $(0, 1)$ & Adecuada para salidas de probabilidad \\
Tanh & $g(z) = \frac{e^z - e^{-z}}{e^z + e^{-z}}$ & $(-1, 1)$ & Centrada en cero \\
ReLU & $g(z) = \max(0, z)$ & ${[}0, +\infty)$ & Computacionalmente eficiente, la más usada \\
\bottomrule
\end{tabular}
\caption{Comparación de funciones de activación comunes}
\end{table}

\begin{warningbox}{¿Por qué es indispensable usar una función de activación no lineal?}
Si no se usa una función de activación no lineal (es decir, si se usa una activación lineal $g(z)=z$), entonces, sin importar cuántas capas tenga la red, la salida de toda la red seguirá siendo solo una combinación lineal de la entrada. Esto significa que una red de múltiples capas se degrada a un modelo lineal de una sola capa, incapaz de aprender ningún patrón no lineal. Solo al introducir no linealidad la red puede aproximar funciones arbitrariamente complejas.
\end{warningbox}

\begin{figure}[H]
\centering
\includegraphics[width=0.85\textwidth]{slides-images/slide-028.jpg}
\caption{Las funciones de activación no lineales permiten que la red neuronal aprenda fronteras de decisión no lineales complejas\protect\footnotemark}
\end{figure}
\footnotetext{Diapositivas, página 28. Izquierda: la activación lineal solo puede producir fronteras de decisión lineales; derecha: la activación no lineal puede aproximar funciones arbitrariamente complejas.}

\subsection{Intuición geométrica del perceptrón}

Amini construye la intuición mediante un ejemplo bidimensional concreto. Considérese un perceptrón con dos entradas $(x_1, x_2)$, con parámetros $w_0=1$, $\mathbf{W}=[3, -2]^T$:

$$\hat{y} = g(1 + 3x_1 - 2x_2)$$

\begin{figure}[H]
\centering
\includegraphics[width=0.85\textwidth]{slides-images/slide-032.jpg}
\caption{Ejemplo de perceptrón: la frontera de decisión divide el espacio bidimensional en dos regiones\protect\footnotemark}
\end{figure}
\footnotetext{Diapositivas, página 32.}

La parte lineal dentro del paréntesis, $z = 1 + 3x_1 - 2x_2$, define una recta en el espacio bidimensional. Esta recta es la \textbf{frontera de decisión}:
\begin{itemize}
\item Cuando $z > 0$, después de la activación Sigmoid $\hat{y} > 0.5$, se predice la clase positiva
\item Cuando $z < 0$, $\hat{y} < 0.5$, se predice la clase negativa
\item Cuando $z = 0$, $\hat{y} = 0.5$, exactamente sobre la frontera de decisión
\end{itemize}

Tomando como ejemplo la entrada $\mathbf{X}=[-1, 2]^T$: $\hat{y} = g(1 + 3\times(-1) - 2\times 2) = g(-6) \approx 0.002$, muy por debajo de 0.5, por lo que este punto se clasifica como clase negativa.

\begin{knowledgebox}{Limitación de un solo perceptrón}
Un solo perceptrón solo puede aprender fronteras de decisión lineales (una recta en el espacio bidimensional, un hiperplano en espacios de mayor dimensión). Esto significa que no puede resolver problemas no separables linealmente como el XOR. Para tratar patrones más complejos, es necesario combinar múltiples perceptrones en una red de varias capas.
\end{knowledgebox}

\subsection{Resumen de la sección}

El perceptrón es la unidad mínima de la red neuronal; su propagación hacia adelante consta de tres pasos: suma ponderada, suma del sesgo y paso por la función de activación. El término de sesgo permite que la frontera de decisión no tenga que pasar por el origen, y la función de activación no lineal permite que la red aprenda patrones complejos. Un solo perceptrón solo puede representar un clasificador lineal.
\section{Construcción de redes neuronales con perceptrones}

\subsection{Perceptrones de salida múltiple y capas Dense}

Un solo perceptrón tiene una sola salida. Al colocar varios perceptrones en paralelo, con todas las entradas conectadas simultáneamente a cada perceptrón, se forma una \textbf{capa Dense} (capa totalmente conectada).

\begin{figure}[H]
\centering
\includegraphics[width=0.85\textwidth]{slides-images/slide-036.jpg}
\caption{Perceptrón de salida múltiple: todas las entradas conectadas densamente a todas las salidas, es decir, una capa Dense\protect\footnotemark}
\end{figure}
\footnotetext{Diapositiva 36 de las láminas.}

Para la $i$-ésima neurona de salida:
$$z_i = w_{0,i} + \sum_{j=1}^{m} x_j \cdot w_{j,i}$$

Debido a que todas las entradas están conectadas a todas las salidas, este tipo de capa se llama capa \textbf{Dense} (densamente conectada) o \textbf{capa totalmente conectada} (Fully Connected Layer).

Implementar una capa Dense en código es muy sencillo:

\begin{lstlisting}[caption={Implementación de la capa Dense en TensorFlow y PyTorch}]
# TensorFlow
import tensorflow as tf
layer = tf.keras.layers.Dense(units=2)

# PyTorch
import torch.nn as nn
layer = nn.Linear(in_features=m, out_features=2)
\end{lstlisting}

\begin{figure}[H]
\centering
\includegraphics[width=0.95\textwidth]{slides-images/slide-037.jpg}
\caption{Implementación de la capa Dense desde cero: comparación de código entre TensorFlow (izquierda) y PyTorch (derecha)\protect\footnotemark}
\end{figure}
\footnotetext{Diapositiva 37 de las láminas.}

\subsection{Red neuronal de una sola capa oculta}

Al agregar una capa Dense entre la capa de entrada y la capa de salida como \textbf{capa oculta} (Hidden Layer), se forma una red neuronal de una sola capa oculta:

\begin{figure}[H]
\centering
\includegraphics[width=0.85\textwidth]{slides-images/slide-039.jpg}
\caption{Estructura de una red neuronal de una sola capa oculta\protect\footnotemark}
\end{figure}
\footnotetext{Diapositiva 39 de las láminas.}

Cálculo de la capa oculta:
$$z_i = w_{0,i}^{(1)} + \sum_{j=1}^{m} x_j \cdot w_{j,i}^{(1)}$$

Cálculo de la capa de salida (usando como entrada los valores de activación de la capa oculta):
$$\hat{y}_i = g\left(w_{0,i}^{(2)} + \sum_{j=1}^{d_1} g(z_j) \cdot w_{j,i}^{(2)}\right)$$

Donde los superíndices $(1)$ y $(2)$ indican, respectivamente, los parámetros de la primera y la segunda capa, y $d_1$ es el número de neuronas de la capa oculta.

\begin{importantbox}{El significado de la capa oculta}
La capa oculta se llama «oculta» porque, en los datos de entrenamiento, solo podemos observar directamente la entrada y la salida, mientras que la representación de la capa oculta es aprendida automáticamente por la red y resulta «invisible» para un observador externo. Cada neurona de la capa oculta aprende alguna transformación no lineal de la entrada, y estas transformaciones en conjunto aportan una representación de características enriquecida para la predicción final.
\end{importantbox}

\subsection{Redes neuronales profundas}

Al apilar más capas ocultas, se forma una \textbf{red neuronal profunda} (Deep Neural Network). De ahí proviene precisamente el «profunda» en «aprendizaje profundo».

\begin{figure}[H]
\centering
\includegraphics[width=0.85\textwidth]{slides-images/slide-042.jpg}
\caption{Red neuronal profunda: apilamiento de múltiples capas ocultas\protect\footnotemark}
\end{figure}
\footnotetext{Diapositiva 42 de las láminas.}

La fórmula de cálculo para la $i$-ésima neurona de la capa $k$ es:
$$z_{k,i} = w_{0,i}^{(k)} + \sum_{j=1}^{n_{k-1}} g(z_{k-1,j}) \cdot w_{j,i}^{(k)}$$

Donde $n_{k-1}$ es el número de neuronas de la capa $k-1$.

Construir una red profunda en código resulta igual de directo:

\begin{lstlisting}[caption={Redes profundas en TensorFlow y PyTorch}]
# TensorFlow
model = tf.keras.Sequential([
    tf.keras.layers.Dense(n1),
    tf.keras.layers.Dense(n2),
    # ...
    tf.keras.layers.Dense(2)
])

# PyTorch
model = nn.Sequential(
    nn.Linear(m, n1), nn.ReLU(),
    nn.Linear(n1, n2), nn.ReLU(),
    # ...
    nn.Linear(nK, 2)
)
\end{lstlisting}

\subsection{Resumen de la sección}

Partiendo de un solo perceptrón, al colocar varios perceptrones en paralelo se forma una capa Dense, y al apilar varias capas Dense se forma una red neuronal profunda. La salida de cada capa sirve de entrada a la siguiente, extrayendo capa por capa características cada vez más abstractas. La «profundidad» de una red neuronal profunda proviene de esta estructura multinivel.

\section{Aplicación de redes neuronales: cuantificación de la pérdida}

\subsection{Problema de ejemplo: ¿aprobará el examen?}

Amini usa un ejemplo intuitivo para ilustrar el flujo de aplicación de las redes neuronales. El problema es: a partir del número de asistencias del estudiante $x_1$ y las horas dedicadas al proyecto final $x_2$, predecir si el estudiante aprobará el curso 6.S191.

\begin{figure}[H]
\centering
\includegraphics[width=0.85\textwidth]{slides-images/slide-047.jpg}
\caption{Problema de ejemplo: dado el nuevo punto de datos $[4,5]$, ¿se predice que aprobará?\protect\footnotemark}
\end{figure}
\footnotetext{Diapositivas, página 47.}

Al introducir el punto de datos $x^{(1)} = [4, 5]$ en una red con una sola capa oculta, supongamos que la red (con pesos inicializados aleatoriamente) produce $\hat{y}_1 = 0.1$. Pero el estudiante en realidad aprobó el curso (etiqueta real $y^{(1)} = 1$). Esto indica que los pesos actuales de la red aún no son suficientemente buenos: necesitamos una forma de cuantificar la diferencia entre la predicción y el valor real.

\subsection{Función de pérdida}

La \textbf{función de pérdida} (Loss Function) mide la diferencia entre el valor predicho por la red y el valor real.

Para la pérdida de un solo punto de datos:
$$\mathcal{L}\big(f(x^{(i)}; \mathbf{W}), y^{(i)}\big)$$

\begin{figure}[H]
\centering
\includegraphics[width=0.85\textwidth]{slides-images/slide-050.jpg}
\caption{Función de pérdida: mide la desviación entre la predicción y el valor real\protect\footnotemark}
\end{figure}
\footnotetext{Diapositivas, página 50.}

La \textbf{pérdida empírica} (Empirical Loss) sobre todo el conjunto de datos es el promedio de las pérdidas de todas las muestras:
$$J(\mathbf{W}) = \frac{1}{n}\sum_{i=1}^{n}\mathcal{L}\big(f(x^{(i)}; \mathbf{W}), y^{(i)}\big)$$

La pérdida empírica también se conoce como \textbf{función objetivo} (Objective Function), \textbf{función de costo} (Cost Function) o \textbf{riesgo empírico} (Empirical Risk).

\begin{importantbox}{Dos funciones de pérdida esenciales}
\textbf{1. Pérdida de entropía cruzada binaria} (Binary Cross Entropy Loss): se usa en problemas de clasificación (la salida es una probabilidad):
$$J(\mathbf{W}) = -\frac{1}{n}\sum_{i=1}^{n}\left[y^{(i)}\log f(x^{(i)};\mathbf{W}) + (1-y^{(i)})\log(1-f(x^{(i)};\mathbf{W}))\right]$$

\textbf{2. Pérdida de error cuadrático medio} (Mean Squared Error Loss): se usa en problemas de regresión (la salida es un valor continuo):
$$J(\mathbf{W}) = \frac{1}{n}\sum_{i=1}^{n}\left(y^{(i)} - f(x^{(i)};\mathbf{W})\right)^2$$
\end{importantbox}

\begin{figure}[H]
\centering
\begin{subfigure}[t]{0.48\textwidth}
\includegraphics[width=\textwidth]{slides-images/slide-052.jpg}
\caption{Pérdida de entropía cruzada (problema de clasificación)}
\end{subfigure}
\hfill
\begin{subfigure}[t]{0.48\textwidth}
\includegraphics[width=\textwidth]{slides-images/slide-053.jpg}
\caption{Error cuadrático medio (problema de regresión)}
\end{subfigure}
\caption{Dos funciones de pérdida esenciales\protect\footnotemark}
\end{figure}
\footnotetext{Diapositivas, páginas 52--53.}

\begin{warningbox}{Elegir la función de pérdida correcta}
La elección de la función de pérdida debe corresponder con el tipo de tarea:
\begin{itemize}
\item Las tareas de clasificación (aprobar/reprobar, gato/perro) usan la pérdida de entropía cruzada
\item Las tareas de regresión (predecir la calificación de un examen, predecir el precio de una casa) usan el error cuadrático medio
\end{itemize}
Una función de pérdida incorrecta impide que el modelo aprenda correctamente. Por ejemplo, usar el MSE en una tarea de clasificación provoca un entrenamiento lento porque la función de pérdida es poco sensible a los límites de probabilidad.
\end{warningbox}

\subsection{Resumen de la sección}

La función de pérdida es un componente central del entrenamiento de redes neuronales: cuantifica la diferencia entre la predicción del modelo y la etiqueta real. La pérdida empírica es el promedio de las pérdidas sobre todas las muestras de entrenamiento. Las tareas de clasificación suelen usar la pérdida de entropía cruzada, y las tareas de regresión, el error cuadrático medio. El objetivo del entrenamiento es encontrar los parámetros de peso que minimicen la pérdida empírica.
\section{Entrenamiento de redes neuronales}

\subsection{El objetivo de la optimización de la pérdida}

El objetivo de entrenar una red neuronal es encontrar los pesos $\mathbf{W}^*$ que minimizan la pérdida empírica $J(\mathbf{W})$:

$$\mathbf{W}^* = \arg\min_{\mathbf{W}} J(\mathbf{W}) = \arg\min_{\mathbf{W}} \frac{1}{n}\sum_{i=1}^{n}\mathcal{L}\big(f(x^{(i)}; \mathbf{W}), y^{(i)}\big)$$

Cabe señalar que $\mathbf{W}$ incluye todos los pesos y sesgos de \textbf{todas las capas} de la red: $\mathbf{W} = \{\mathbf{W}^{(0)}, \mathbf{W}^{(1)}, \ldots\}$.

\begin{figure}[H]
\centering
\includegraphics[width=0.85\textwidth]{slides-images/slide-057.jpg}
\caption{El landscape de la función de pérdida: el objetivo es encontrar el punto más bajo\protect\footnotemark}
\end{figure}
\footnotetext{Diapositivas, página 57.}

\subsection{Descenso de gradiente}

El \textbf{descenso de gradiente} (Gradient Descent) es el algoritmo central para optimizar la función de pérdida. Su idea básica es sencilla y elegante:

\begin{importantbox}{El algoritmo de descenso de gradiente}
\begin{enumerate}
\item \textbf{Inicializar aleatoriamente} los pesos $\mathbf{W}$
\item \textbf{Calcular el gradiente}: $\frac{\partial J(\mathbf{W})}{\partial \mathbf{W}}$, el cual apunta en la dirección en la que la pérdida crece más rápido
\item \textbf{Actualizar los pesos}: dar un paso pequeño en la dirección \textbf{opuesta} al gradiente
$$\mathbf{W} \leftarrow \mathbf{W} - \eta \frac{\partial J(\mathbf{W})}{\partial \mathbf{W}}$$
\item Repetir los pasos 2--3 hasta la convergencia
\end{enumerate}
Aquí $\eta$ es la \textbf{tasa de aprendizaje} (learning rate), que controla la magnitud de cada actualización.
\end{importantbox}

\begin{figure}[H]
\centering
\includegraphics[width=0.85\textwidth]{slides-images/slide-060.jpg}
\caption{Visualización del descenso de gradiente: un descenso iterativo en la dirección opuesta al gradiente\protect\footnotemark}
\end{figure}
\footnotetext{Diapositivas, página 60.}

Una intuición del gradiente: imagina que estás de pie en una zona de colinas, con los ojos vendados, y quieres llegar al punto más bajo. Lo único que puedes hacer es sentir la pendiente del terreno bajo tus pies (el gradiente) y dar un paso en la dirección de mayor descenso. Al repetir este proceso, terminarás por llegar a alguna depresión del terreno.

\subsection{Retropropagación}

Entonces, ¿cómo calcular de manera eficiente el gradiente $\frac{\partial J}{\partial \mathbf{W}}$? La respuesta es la \textbf{retropropagación} (Backpropagation).

La retropropagación utiliza la \textbf{regla de la cadena} (Chain Rule) para propagar el gradiente hacia atrás, capa por capa, comenzando desde la capa de salida. Para cada peso $w$ de la red, la regla de la cadena permite descomponer el gradiente de la pérdida final respecto a ese peso en el producto de varios gradientes intermedios.

Tomando como ejemplo una red simple de dos capas, el gradiente de la pérdida $J$ respecto al peso $w_1$ de la primera capa puede expresarse como:

$$\frac{\partial J}{\partial w_1} = \frac{\partial J}{\partial \hat{y}} \cdot \frac{\partial \hat{y}}{\partial z_2} \cdot \frac{\partial z_2}{\partial z_1} \cdot \frac{\partial z_1}{\partial w_1}$$

\begin{importantbox}{El núcleo de la retropropagación}
La esencia de la retropropagación es la \textbf{aplicación recursiva de la regla de la cadena}. Se calculan gradientes locales empezando por el extremo de salida y se propagan hacia atrás, capa por capa, multiplicando todos los gradientes locales para obtener el gradiente final. Esto hace posible calcular gradientes de manera eficiente en redes profundas con millones de parámetros.
\end{importantbox}

Los marcos de aprendizaje profundo modernos (TensorFlow, PyTorch) realizan automáticamente el cálculo del gradiente de la retropropagación mediante la \textbf{diferenciación automática} (Automatic Differentiation), sin que el desarrollador tenga que derivarlo a mano.

\begin{lstlisting}[caption={Diferenciación automática en TensorFlow y PyTorch}]
# TensorFlow: GradientTape
import tensorflow as tf
with tf.GradientTape() as tape:
    loss = compute_loss(model, x, y)
grads = tape.gradient(loss, model.trainable_variables)

# PyTorch: autograd
import torch
loss = compute_loss(model, x, y)
loss.backward()
# Gradients stored in param.grad
\end{lstlisting}

\subsection{La importancia de la tasa de aprendizaje}

La tasa de aprendizaje $\eta$ es uno de los hiperparámetros más determinantes al entrenar una red neuronal. Su valor incide directamente en el éxito o el fracaso del entrenamiento:

\begin{table}[H]
\centering
\begin{tabular}{lp{9cm}}
\toprule
\textbf{Tasa de aprendizaje} & \textbf{Efecto} \\
\midrule
Muy baja & La convergencia es extremadamente lenta y puede quedar atrapada en un mínimo local \\
Muy alta & El entrenamiento es inestable, la pérdida oscila repetidamente cerca del mínimo o incluso diverge \\
Adecuada & Converge de manera estable a un buen mínimo \\
\bottomrule
\end{tabular}
\caption{Efecto del valor de la tasa de aprendizaje sobre el entrenamiento}
\end{table}

\begin{warningbox}{Trampas comunes al ajustar la tasa de aprendizaje}
\begin{itemize}
\item Una tasa de aprendizaje demasiado alta no solo provoca oscilaciones, sino que puede hacer que la pérdida crezca cada vez más (diverja), porque cada actualización rebasa el mínimo
\item Una tasa de aprendizaje demasiado baja es segura, pero puede requerir un tiempo de entrenamiento sumamente largo y con facilidad queda atrapada en mínimos locales poco profundos
\item En la práctica se usan con frecuencia métodos de \textbf{tasa de aprendizaje adaptativa} (como Adam, RMSProp), que ajustan automáticamente la tasa de aprendizaje de cada parámetro a partir del historial de sus gradientes
\end{itemize}
\end{warningbox}

\begin{importantbox}{Algoritmos de tasa de aprendizaje adaptativa}
En el entrenamiento real rara vez se usa una tasa de aprendizaje fija. Entre los optimizadores adaptativos más comunes están:
\begin{itemize}
\item \textbf{SGD with Momentum}: usa un promedio móvil exponencial del gradiente para acelerar la convergencia
\item \textbf{Adam} (Adaptive Moment Estimation): combina estimaciones del primer y del segundo momento para ajustar la tasa de aprendizaje de cada parámetro por separado
\item \textbf{RMSProp}: normaliza la tasa de aprendizaje mediante un promedio móvil del cuadrado del gradiente
\end{itemize}
La idea central de estos algoritmos es que, en lugar de usar una misma tasa de aprendizaje fija para todos los parámetros, esta se ajusta de manera dinámica según el historial del gradiente de cada parámetro.
\end{importantbox}

\subsection{Descenso de gradiente por minilotes}

El descenso de gradiente estándar necesita calcular la pérdida y el gradiente sobre \textbf{todo el conjunto de datos} antes de poder realizar una sola actualización de los pesos, lo cual implica un costo de cómputo enorme en conjuntos de datos a gran escala.

El \textbf{descenso de gradiente estocástico} (Stochastic Gradient Descent, SGD) representa el otro extremo: en cada paso se usa \textbf{una sola muestra} para calcular el gradiente y actualizar los pesos; es rápido, pero muy ruidoso.

La solución intermedia que se usa en la práctica es el \textbf{descenso de gradiente por minilotes}: en cada paso se toma al azar un lote pequeño (minilote) de muestras del conjunto de datos para calcular el gradiente:

$$\frac{\partial J(\mathbf{W})}{\partial \mathbf{W}} \approx \frac{1}{B}\sum_{k=1}^{B}\frac{\partial \mathcal{L}(f(x^{(k)}; \mathbf{W}), y^{(k)})}{\partial \mathbf{W}}$$

Aquí $B$ es el tamaño de lote.

\begin{table}[H]
\centering
\begin{tabular}{lp{4.5cm}p{4.5cm}}
\toprule
\textbf{Método} & \textbf{Ventajas} & \textbf{Desventajas} \\
\midrule
GD por lote completo & Gradiente exacto, convergencia estable & Alto costo de cómputo y gran demanda de memoria \\
SGD ($B=1$) & Frecuencia de actualización muy alta & Gradiente muy ruidoso, convergencia inestable \\
SGD por minilotes & Equilibra velocidad y estabilidad, se puede paralelizar & Requiere ajustar el tamaño de lote \\
\bottomrule
\end{tabular}
\caption{Comparación de tres estrategias de descenso de gradiente}
\end{table}

\begin{knowledgebox}{Beneficios adicionales de los minilotes}
El descenso de gradiente por minilotes no es solo una solución intermedia en términos de eficiencia computacional; el ruido de gradiente que introduce en realidad ayuda al modelo a escapar de mínimos locales pronunciados y tiende a converger hacia regiones de mínimos más planas, las cuales suelen tener mejor capacidad de generalización. Además, el hardware de las GPU está naturalmente diseñado para el cómputo paralelo sobre datos en lote, lo que hace que la velocidad real de entrenamiento del SGD por minilotes supere ampliamente a la actualización muestra por muestra.
\end{knowledgebox}

\subsection{Los retos del descenso de gradiente: mínimos locales y puntos de silla}

En espacios de parámetros de alta dimensión, el landscape de la función de pérdida es sumamente complejo. El descenso de gradiente enfrenta varios retos clave:

\begin{enumerate}
\item \textbf{Mínimos locales} (Local Minima): la función de pérdida puede tener muchos mínimos locales, y el descenso de gradiente puede converger a una posición que no es el óptimo global
\item \textbf{Puntos de silla} (Saddle Points): en espacios de alta dimensión, los puntos de silla son más comunes que los mínimos locales. En un punto de silla, el gradiente es cero en algunas direcciones, lo que estanca la optimización
\item \textbf{Mesetas} (Plateaus): regiones donde el gradiente de la función de pérdida es cercano a cero, lo que hace que las actualizaciones sean extremadamente lentas
\end{enumerate}

\begin{knowledgebox}{Intuición sobre la optimización en espacios de alta dimensión}
En espacios de baja dimensión (como dos dimensiones), los mínimos locales parecen muy comunes. Pero en los espacios de parámetros de millones de dimensiones típicos de las redes neuronales, para que se forme un verdadero mínimo local se necesita que \textbf{todas las dimensiones} se curven hacia arriba al mismo tiempo, lo cual es estadísticamente muy improbable. En realidad, el problema más común en la optimización de alta dimensión son los puntos de silla, en los que algunas dimensiones suben y otras bajan. Esto también explica por qué el aprendizaje profundo funciona tan bien en la práctica: aunque el landscape de la pérdida parezca complejo, la mayoría de los «pozos» tienen una vía de escape.
\end{knowledgebox}

\subsection{Resumen de la sección}

El núcleo del entrenamiento de una red neuronal consiste en minimizar la función de pérdida mediante el descenso de gradiente. La retropropagación utiliza la regla de la cadena para calcular el gradiente de manera eficiente. La tasa de aprendizaje es el hiperparámetro más determinante: tanto un valor demasiado alto como uno demasiado bajo generan problemas, y los algoritmos de tasa de aprendizaje adaptativa (como Adam) pueden ajustarla automáticamente. El SGD por minilotes logra un buen equilibrio entre eficiencia computacional y calidad del gradiente. Aunque la optimización en espacios de alta dimensión enfrenta retos como los mínimos locales y los puntos de silla, los algoritmos de optimización modernos tienen un desempeño sobresaliente en la práctica.
\section{Regularización y generalización}

\subsection{El problema del sobreajuste}

Al entrenar redes neuronales, un desafío central es el \textbf{sobreajuste} (Overfitting): el modelo tiene un buen desempeño en los datos de entrenamiento, pero un desempeño deficiente en datos de prueba que nunca ha visto. Esto significa que el modelo "memorizó" el ruido y los patrones particulares de los datos de entrenamiento, sin aprender reglas verdaderamente generalizables.

\begin{importantbox}{Subajuste, ajuste adecuado y sobreajuste}
\begin{itemize}
\item \textbf{Subajuste} (Underfitting): el modelo es demasiado simple para capturar los patrones fundamentales de los datos y tiene un desempeño deficiente tanto en el conjunto de entrenamiento como en el de prueba
\item \textbf{Ajuste adecuado} (Good Fit): el modelo aprende de manera apropiada los patrones centrales de los datos y tiene buen desempeño tanto en el conjunto de entrenamiento como en el de prueba
\item \textbf{Sobreajuste} (Overfitting): el modelo es demasiado complejo y se ajusta al ruido de los datos de entrenamiento; tiene buen desempeño en el conjunto de entrenamiento pero deficiente en el de prueba
\end{itemize}
\end{importantbox}

\subsection{Técnicas de regularización}

Para mitigar el sobreajuste, es necesario usar técnicas de \textbf{regularización} (Regularization). Amini presentó dos métodos principales:

\textbf{1. Dropout}

La idea central de Dropout es sumamente sencilla: durante el entrenamiento, con probabilidad $p$ se "apaga" (se pone en cero) de manera aleatoria una parte de las neuronas.

\begin{itemize}
\item En la propagación hacia adelante de cada mini-batch, se elige de manera aleatoria alrededor de $p\%$ de las neuronas para que su salida sea cero
\item Esto obliga a la red a no depender de ninguna neurona individual, sino a aprender representaciones de características más distribuidas y robustas
\item Durante la prueba no se usa Dropout, pero es necesario escalar la salida de manera proporcional
\end{itemize}

\begin{lstlisting}[caption={Implementación en código de Dropout}]
# TensorFlow
tf.keras.layers.Dropout(rate=0.5)

# PyTorch
nn.Dropout(p=0.5)
\end{lstlisting}

\textbf{2. Early Stopping}

Early Stopping monitorea el desempeño del modelo en el conjunto de validación: cuando la pérdida de validación deja de disminuir (o incluso empieza a aumentar), se detiene el entrenamiento. Esto aprovecha una observación clave: el sobreajuste suele ocurrir en las etapas finales del entrenamiento, momento en el que el modelo empieza a "memorizar" los datos de entrenamiento en lugar de aprender reglas generalizables.

\begin{warningbox}{La regularización no es opcional}
En la práctica, una red neuronal profunda sin regularización casi con certeza se sobreajustará, en especial cuando el número de parámetros del modelo es mucho mayor que el número de ejemplos de entrenamiento (algo muy común en el aprendizaje profundo actual). Dropout y Early Stopping deben ser prácticas estándar, no elementos adicionales opcionales.
\end{warningbox}

\subsection{Resumen de la sección}

El sobreajuste es un desafío central en el entrenamiento del aprendizaje profundo. Dropout obliga a la red a aprender características robustas apagando neuronas de manera aleatoria, y Early Stopping detiene el entrenamiento en el momento adecuado monitoreando la pérdida en el conjunto de validación. La combinación de ambos puede mejorar de manera significativa la capacidad de generalización del modelo.
\section{Síntesis y ampliación}

\subsection{Repaso de los puntos clave de esta clase}

La primera clase de MIT 6.S191 presentó de manera sistemática, de lo general a lo particular, el marco fundamental del aprendizaje profundo:

\begin{enumerate}
\item \textbf{Ubicación del aprendizaje profundo}: IA $\supset$ ML $\supset$ DL, cuyo núcleo consiste en aprender de manera automática representaciones de características a partir de datos crudos
\item \textbf{Perceptrón}: la unidad básica de las redes neuronales, que consta de tres pasos: suma ponderada, sesgo y activación no lineal
\item \textbf{Funciones de activación}: funciones no lineales como Sigmoid, Tanh y ReLU que permiten a la red aproximar funciones arbitrariamente complejas
\item \textbf{Redes profundas}: se construyen apilando múltiples capas Dense, y cada capa aprende características cada vez más abstractas
\item \textbf{Funciones de pérdida}: la entropía cruzada (clasificación) y el MSE (regresión) miden respectivamente la distancia entre la predicción y el valor real
\item \textbf{Proceso de entrenamiento}: descenso de gradiente + retropropagación + una tasa de aprendizaje adecuada
\item \textbf{Regularización}: Dropout y Early Stopping previenen el sobreajuste
\end{enumerate}

\begin{table}[H]
\centering
\begin{tabular}{lp{5cm}p{5cm}}
\toprule
\textbf{Concepto} & \textbf{Expresión matemática} & \textbf{Explicación intuitiva} \\
\midrule
Perceptrón & $\hat{y} = g(w_0 + \mathbf{X}^T\mathbf{W})$ & Voto ponderado + decisión no lineal \\
Capa Dense & $z_i = w_{0,i} + \sum_j x_j w_{j,i}$ & Varios perceptrones en paralelo \\
Función de pérdida & $J(\mathbf{W}) = \frac{1}{n}\sum_i \mathcal{L}$ & Medida cuantitativa de la calidad de la predicción \\
Descenso de gradiente & $\mathbf{W} \leftarrow \mathbf{W} - \eta \nabla J$ & Bajar la montaña en la dirección más pronunciada \\
Retropropagación & Aplicación recursiva de la regla de la cadena & Propaga el error capa por capa, de la salida hacia la entrada \\
Dropout & Probabilidad $p$ de anular aleatoriamente & Fuerza la redundancia y evita la dependencia de una sola característica \\
\bottomrule
\end{tabular}
\caption{Tabla de referencia rápida de los conceptos clave de esta clase}
\end{table}

\subsection{De lo fundamental a lo más avanzado}

Amini enfatizó varias veces durante la clase que el contenido de esta sesión (el perceptrón, el descenso de gradiente, la retropropagación) son invenciones de hace décadas, pero constituyen la base necesaria para entender el aprendizaje profundo moderno. Las siguientes clases del curso cubrirán:

\begin{itemize}
\item \textbf{Modelado de secuencias}: arquitecturas como RNN y LSTM para procesar datos temporales
\item \textbf{Visión por computadora}: CNN y arquitecturas de visión modernas
\item \textbf{Modelos generativos}: VAE, GAN, Diffusion Models
\item \textbf{Aprendizaje por refuerzo}: hacer que el modelo aprenda una política mediante la interacción con el entorno
\item \textbf{Modelos de lenguaje de gran escala}: la arquitectura Transformer y los LLM modernos
\end{itemize}

\subsection{Lecturas adicionales}

\begin{itemize}
\item Sitio oficial del curso: \url{https://introtodeeplearning.com}
\item Repositorio de código de los laboratorios: \url{https://github.com/MITDeepLearning/introtodeeplearning}
\item Ian Goodfellow, Yoshua Bengio, Aaron Courville, \textit{Deep Learning}, MIT Press, 2016
\item Michael Nielsen, \textit{Neural Networks and Deep Learning}, libro de texto gratuito en línea
\item 3Blue1Brown, serie de videos \textit{Neural Networks} (visualización intuitiva)
\end{itemize}

%% --- Fin del contenido --- %%

\end{document}
